\lyseis{5}{Χαρακτήρες διαιρετότητας: ξέρω χωρίς να διαιρέσω}
\sol{1} Με το 2: 3.140, 1.234, 5.050. Με το 5: 3.140, 2.025, 5.050. Με το 10: 3.140, 5.050.
\sol{2} Άθροισμα ψηφίων: 123 $\to$ 6 (ναι), 4.521 $\to$ 12 (ναι), 1.001 $\to$ 2 (όχι), 99.999 $\to$ 45 (ναι), 2.468 $\to$ 20 (όχι). Έλεγχος, π.χ. $4.521:3=1.507$.
\sol{3} 1.089 $\to$ 18 (ναι), 4.527 $\to$ 18 (ναι), 81.081 $\to$ 18 (ναι), 7.777 $\to$ 28 (όχι), 90.009 $\to$ 18 (ναι).
\sol{4} 316 (16, ναι), 1.002 (02, όχι), 7.500 (00, ναι), 2.024 (24, ναι), 98 (98 δεν διαιρείται με το 4, όχι).
\sol{5} 175 (75, ναι), 2.350 (50, ναι), 4.600 (00, ναι), 1.065 (65, όχι), 3.025 (25, ναι).
\sol{6} \begin{tabular}[t]{@{}lccccccc@{}}\toprule & 2 & 3 & 4 & 5 & 9 & 10 & 25\\\midrule 840 & ναι & ναι & ναι & ναι & όχι & ναι & όχι\\ 1.575 & όχι & ναι & όχι & ναι & ναι & όχι & ναι\\ 6.048 & ναι & ναι & ναι & όχι & ναι & όχι & όχι\\ 20.250 & ναι & ναι & όχι & ναι & ναι & ναι & ναι\\\bottomrule\end{tabular}
\sol{7} Το 54 (άρτιος και $5+4=9$), το 4.230 (άρτιος και $4+2+3+0=9$) και το 7.776 (άρτιος και $7+7+7+6=27$). Το 128 δεν διαιρείται με το 3, και το 1.005 είναι περιττός.
\sol{8} α) $4+7+2=13$, χρειάζεται το 18, άρα $\square=5$: 4.752. \quad β) $3+5=8$ και πρέπει να φτάσει σε πολλαπλάσιο του 3: $\square=1,\,4,7$ (315, 345, 375). \quad γ) Τα δύο τελευταία ψηφία $\square 4$ πρέπει να διαιρούνται με το 4: $04, 24, 44, 64, 84$, άρα $\square=0,\,2,4,\,6,8$. \quad δ) Πρέπει να τελειώνει σε 0 (διαιρείται με 2 και 5): 170.
\sol{9} α) Λάθος: το 14 τελειώνει σε 4 και δεν διαιρείται. β) Σωστό: το 9 διαιρείται με το 3. γ) Σωστό. δ) Λάθος: το 12 διαιρείται και με τα δύο, αλλά όχι με το 24. ε) Σωστό. στ) Σωστό.
\sol{10} Το 3.465 έχει άθροισμα ψηφίων 18, άρα διαιρείται με το 3 και το 9. Τελειώνει σε 5, άρα διαιρείται με το 5. Τα δύο τελευταία ψηφία 65 δεν διαιρούνται με το 4, και δεν τελειώνει σε 0 (όχι 4, όχι 10). Σε 3 ομάδες: $3.465:3=1.155$ €. Σε 5: $693$ €. Σε 9: $385$ €.
\sol{11} Το 140 τελειώνει σε 0, άρα διαιρείται με το 2, το 5 και το 10. Τα δύο τελευταία ψηφία 40 διαιρούνται με το 4, άρα και με το 4. Το άθροισμα ψηφίων είναι 5, άρα όχι με το 3, και επομένως ούτε με το 6. Γίνεται με 2, 4, 5 και 10 στοίβες.
\sol{12} Ο αριθμός είναι $72\square 0$, με άθροισμα ψηφίων $9+\square$. Πρέπει να διαιρείται με το 9, άρα $\square=0$ ή $9$. Επειδή το τρίτο ψηφίο είναι μεγαλύτερο από το 2, είναι το 9: ο κωδικός είναι 7.290.
\sol{13} Με το 5: $60:5=12$ θέσεις. Με το 10: 6 θέσεις (10, 20, \dots, 60). Και με το 3 και με το 5: οι 15, 30, 45, 60, δηλαδή 4 θέσεις.
\sol{14} Το 2.028 είναι δίσεκτο ($28:4=7$). Το 2.000 είναι δίσεκτο (διαιρείται με το 400). Το 1.900 δεν είναι (διαιρείται με το 100 αλλά όχι με το 400). Το 2.100 δεν είναι. Το 2.026 δεν είναι (26 δεν διαιρείται με το 4).
\sol{15} Αν ένας αριθμός τελειώνει σε 0, είναι πολλαπλάσιο του 10: $10\cdot\kappa$. Επειδή $10=2\cdot 5$, έχουμε $10\kappa=2\cdot(5\kappa)=5\cdot(2\kappa)$. Άρα διαιρείται και με το 2 και με το 5.
\sol{16} Τα κομμάτια $1\cdot 999$, $4\cdot 99$ και $5\cdot 9$ διαιρούνται με το 9 και με το 3. Αρκεί να ελέγξουμε το $1+4+5+2=12$. Το 12 διαιρείται με το 3 (το 1.452 διαιρείται με το 3) αλλά όχι με το 9 (το 1.452 δεν διαιρείται με το 9).
\sol{17} Πρέπει να διαιρείται με το 4, το 9 και το 5, και τότε διαιρείται και με το 2 και με το 3. Τα πολλαπλάσια του $4\cdot 9\cdot 5=180$ είναι 180, 360, \dots Το πρώτο τετραψήφιο είναι $180\cdot 6=1.080$.
\sol{18} Το μεσαίο ψηφίο είναι το άθροισμα των δύο άκρων: $1+1=2$, $1+2=3$, $1+3=4$, $1+4=5$. Το 165 ταιριάζει ($1+5=6$) και $165=11\cdot 15$. Το 175 δεν ταιριάζει ($1+5=6\neq 7$) και $175:11$ δίνει υπόλοιπο 10.
\sol{19} Για το 3 μετράνε όλα τα ψηφία, όχι μόνο το τελευταίο. Το άθροισμα είναι $2+0+2+6=10$, που δεν διαιρείται με το 3, άρα ο 2.026 δεν διαιρείται (η διαίρεση δίνει $2.026=3\cdot 675+1$).
